\chapter{Bevezetés}
\label{ch:intro}

A Scalable Vector Graphics (SVG)~\cite{w3c-svg} XML-alapú leírónyelv kétdimenziós vektorgrafikák készítéséhez. Első változatát a W3C 2001-ben tette ajánlássá, hogy a webnek legyen szabványos vektorgrafikus formátuma~\cite{w3c-svg10}. Az SVG a képet alakzatok, szövegek és beágyazott képek matematikai leírásaként tárolja, így az minőségromlás nélkül tetszőlegesen nagyítható. A rasztergrafika~\cite{raster-graphics} ezzel szemben képpontok rácsával ábrázolja a képet, és inkább fényképekhez alkalmas. Az SVG mára a web egyik alapformátuma. A W3Techs 2026. novemberi felmérése szerint a weboldalak 68,8\%-a használ SVG-képet~\cite{w3techs-svg}, jellemzően ikonokhoz, logókhoz és diagramokhoz.

\medskip
Egy SVG-fájl szövegszerkesztővel is módosítható, de egy görbe vagy egy összetett alakzat így nehezen kezelhető, ezért a gyakorlatban grafikus szerkesztőt használnak hozzá. \Aref{ch:svg}. fejezetben vizsgált szerkesztők között van nagy tudású, de asztali vagy fizetős program, és van egyszerű, de korlátozott webes eszköz is. Szakdolgozatom célja egy olyan webalapú SVG-szerkesztő tervezése és megvalósítása, amely a böngészőben, telepítés és regisztráció nélkül fut, mégis teljes értékű útvonalszerkesztést kínál, és telefonon is kezelhető. Az elkészült alkalmazás a \textit{Solid SVG} nevet kapta. A feladatkiírás külön kérdése, hogy a beolvasott SVG-fájlok milyen esetekben és milyen korlátozásokkal módosíthatók, hiszen tartalmazhatnak szkripteket vagy a szerkesztő által nem kezelt szerkezeteket.

\medskip
Egy ilyen szerkesztő elsőre egyszerű feladatnak tűnhet, hiszen a böngésző maga is meg tudja jeleníteni az SVG ábrákat. A nehézségek a szerkesztésnél kezdődnek. Egy elforgatott alakzat átméretezésekor a szemközti saroknak a helyén kell maradnia, két alakzat egyesítésekor a görbéknek görbének kell maradniuk, az útvonallá alakított szövegnek pedig pontosan úgy kell kinéznie, mint a képernyőn, akkor is, ha a betűtípus egy másik gépen nincs telepítve. Mindezt szerver nélkül kell megoldani, úgy, hogy minden művelet visszavonható legyen. A megvalósításhoz a Svelte keretrendszert~\cite{svelte-overview} és a TypeScript nyelvet~\cite{mdn-ts} választottam.

\medskip
\Aref{ch:svg}. fejezet az SVG grafikus elemeit és a hasonló célú alkalmazásokat mutatja be, és ezek alapján megfogalmazza a célkitűzéseket. \Aref{ch:requirements}. fejezet a követelményeket és az alkalmazás tervét ismerteti, \aref{ch:implementation}. fejezet a megvalósítás érdekesebb részeit, \aref{ch:test}. fejezet a tesztelést és a használatot. \Aref{ch:evaluation}. fejezet a célkitűzések teljesülését értékeli, \aref{ch:summary}. fejezet pedig magyarul és angolul is összefoglalja a munkát.
